\documentclass[11pt]{article}

\usepackage{amsmath,amssymb,amsthm,mathtools}
\usepackage[margin=1in]{geometry}

\newtheorem{theorem}{Theorem}

\begin{document}

\title{An Exact Fractional-Fibre Counterexample for Atlas 152}
\author{}
\date{}
\maketitle

\section{Statement and identifying data}

Let $H$ be Atlas graph $152$, with graph6 code \texttt{EhNG}, vertex set
$\{0,1,2,3,4,5\}$, and edge set
\[
 \{01,05,12,15,23,24,34,45\}.
\]
Thus $v(H)=6$ and $e(H)=8$.  The vertices $0,1,5$ span a triangle, so
$\chi(H)\geq3$.  The colour classes
\[
 \{1,4\},\qquad \{2,5\},\qquad \{0,3\}
\]
give a proper three-colouring, and hence $\chi(H)=3$.  Exact
deletion--contraction gives
\begin{equation}
 \chi_H(x)=x(x-1)(x-2)^2(x^2-3x+3).
 \label{eq:chromatic-polynomial}
\end{equation}
Consequently the conjectured lower target is
\begin{equation}
 \Phi_H(p)=p(2p-1)^2(3p^2-3p+1).
 \label{eq:target}
\end{equation}
In particular the admissible range begins at $p=1/2$.

\begin{theorem}
There is a five-step graphon $W$ of edge density
\[
 p=\frac{18750093281}{37500000000}>\frac12
\]
for which
\[
 t(H,W)-\Phi_H(p)
 =-
 \frac{48407747400063658542114896407207170401}
 {6179809570312500000000000000000000000000000000000000}<0.
\]
Therefore Atlas $152$ is a negative case of the proposed inequality.
\end{theorem}

\section{The five-step graphon}

It is useful first to give a one-parameter family.  For
$0<\varepsilon<1/2$, partition the probability space into
\[
 A_1,\ A_2,\ B,\ S_1,\ S_2
\]
with respective masses
\begin{equation}
 \left(\frac12-\varepsilon\right)\frac{624}{625},\qquad
 \left(\frac12-\varepsilon\right)\frac1{625},\qquad
 \frac12,\qquad \frac{\varepsilon}{10},\qquad
 \frac{9\varepsilon}{10}.
 \label{eq:masses}
\end{equation}
Define a symmetric graphon by putting
\begin{equation}
 W_{A_1B}=W_{A_2B}=W_{A_1S_1}=W_{A_2S_2}=W_{BS_2}=1,
 \qquad W_{BS_1}=\frac1{16},
 \label{eq:cells}
\end{equation}
and putting $W=0$ on every other pair of partition classes.  All cell
values lie in $[0,1]$, and the five masses in \eqref{eq:masses} are positive
and sum to one.

This is an arbitrary-fibre perturbation of the balanced complete bipartite
graphon.  Relative to the bulk parts $A=A_1\sqcup A_2$ and $B$, the two
exceptional profiles are
\[
 (u_1,v_1)=\left(\frac{624}{625},\frac1{16}\right),
 \qquad
 (u_2,v_2)=\left(\frac1{625},1\right),
\]
with exceptional profile masses $1/10$ and $9/10$.  Their density
displacements are
\begin{equation}
 d_1=u_1+v_1-1=\frac{609}{10000},
 \qquad d_2=u_2+v_2-1=\frac1{625}.
 \label{eq:displacements}
\end{equation}
Both are strictly positive.

\section{Exact local mechanism}

The $A$-fibres are the complementary indicators of $A_1$ and $A_2$.
The $B$-fibres are the constant functions $1/16$ and $1$.  Therefore their
two Gram matrices are
\[
 (h_A(i,j))=
 \begin{pmatrix}624/625&0\\0&1/625\end{pmatrix},
 \qquad
 (h_B(i,j))=
 \begin{pmatrix}1/256&1/16\\1/16&1\end{pmatrix}.
\]
For Atlas $152$, the exact joint kernel from the arbitrary-fibre threshold
calculation is
\[
 \mathcal A(i,j)
 =v_i h_A(i,j)+u_i h_B(i,j)+u_iv_i.
\]
Here
\[
 (\mathcal A(i,j))=
 \begin{pmatrix}
 1287/10000&78/625\\[2pt]
 17/10000&3/625
 \end{pmatrix}.
\]
At zero exceptional--exceptional cell value, eight times the polarized
quadratic kernel is
\[
 M_{ij}=\mathcal A(i,j)\mathcal A(j,i)-4d_id_j,
\]
so
\begin{equation}
 M=
 \begin{pmatrix}
 34569/20000000&-111/625000\\[2pt]
 -111/625000&1/78125
 \end{pmatrix},
 \qquad
 \det M=-\frac{2943}{312500000000}<0.
 \label{eq:negative-kernel-matrix}
\end{equation}
For the exceptional mass vector $w=(1/10,9/10)$,
\begin{equation}
 w^{\mathsf T}Mw=-\frac{8631}{2000000000}<0.
 \label{eq:negative-quadratic}
\end{equation}
Thus the coefficient of $\varepsilon^2$ in the full gap is
$-8631/16000000000$.  This calculation explains the construction, but the
classification below uses a fixed exact value of $\varepsilon$ and hence
does not rely on an asymptotic remainder.

\section{Exact density and homomorphism density}

Summing the ordered cells in \eqref{eq:cells} gives
\begin{equation}
 p(\varepsilon)
 =\frac12+\frac{753}{100000}\varepsilon
 -\frac{20256}{100000}\varepsilon^2.
 \label{eq:edge-density}
\end{equation}
For transparency, the homomorphism-density calculation can be grouped by
the numbers of vertices sent to
$(A_1,A_2,B,S_1,S_2)$.  All nonzero assignments fall into the following
three groups; the second column is the sum of their edge-cell weights.
\[
\begin{array}{c|c}
\text{class-count vector}&\text{total cell weight}\\ \hline
(2,0,2,2,0)&1089/32768\\
(1,1,2,1,1)&17/32\\
(0,2,2,0,2)&18
\end{array}
\]
Multiplying by the corresponding class masses from \eqref{eq:masses} and
adding yields the compact identity
\begin{equation}
 t(H,W_\varepsilon)
 =\frac{2224881}{80000000000}\,
 \varepsilon^2(1-2\varepsilon)^2.
 \label{eq:hom-density}
\end{equation}
This table is also reconstructed by direct exact enumeration of the
$5^6$ vertex assignments in
\texttt{codes/needle\_higher\_order\_tests.py}.

Now set $\varepsilon=1/3000$.  The five class masses are
\[
 \frac{38974}{78125},\qquad
 \frac{1499}{1875000},\qquad
 \frac12,\qquad
 \frac1{30000},\qquad
 \frac3{10000}.
\]
Equations \eqref{eq:edge-density} and \eqref{eq:hom-density} give
\begin{equation}
 p=\frac{18750093281}{37500000000}
 =\frac12+\frac{93281}{37500000000},
 \qquad
 t(H,W)=\frac{185159623403}{60000000000000000000000}.
 \label{eq:witness-values}
\end{equation}
Substitution of this exact $p$ into \eqref{eq:target} gives
\[
 \Phi_H(p)=
 \frac{19119261293088692808932739896407207170401}
 {6179809570312500000000000000000000000000000000000000}.
\]
Putting the two quantities over the same denominator produces the strictly
negative gap stated in the theorem.

\section{Why the earlier positive subcones do not apply}

The construction lies precisely in the residual cone left by the preceding
Atlas $152$ analysis.
\begin{itemize}
 \item The two mean profiles are incomparable: $u_1>u_2$ but $v_1<v_2$.
 \item Their density displacements in \eqref{eq:displacements} are unequal.
 \item The first profile is strict-interior: $u_1<1$, $v_1<1$, and
       $u_1+v_1>1$.
 \item Although both $A$-fibres are indicators, the first $B$-fibre is the
       genuinely fractional constant $1/16$.  Thus the two-profile theorem
       requiring indicator fibres in both bulk parts does not apply.
 \item The profiles are not both on the axes, are not a coordinatewise
       chain, and are not on one constant-displacement line.
\end{itemize}
Accordingly none of the preceding positive subtheorems applies.  In
particular, the indicator, axis-profile, and comparable-mean results do not
apply.  Neither do the equal-displacement, bounded-dispersion, or comonotone
results.

\section{Hidden-assumption audit and limitations}

The witness is a finite-step graphon on a genuine probability space; no
finite-graph approximation, injective-copy interpretation, regularity,
minimum-degree condition, or limiting argument is used.  Every cell value
is in $[0,1]$, every class has positive rational mass, and the edge density
is in the complete required range $p\geq1/2$.  Both $p$ and the
homomorphism density use ordered graphon integrals.  Formula
\eqref{eq:target} is the polynomial continuation of the conjectured target,
but the witness has $p<1$, so no endpoint convention is needed.

This construction classifies only Atlas $152$.  It does not contradict the
universal high-density theorem: its edge density is just above $1/2$, far
outside the small neighbourhood of $p=1$ certified there.  Nor does it
produce a general closure operation or a counterexample for any other
currently open row.

\appendix
\section{What the Lean formalization actually proves}

The theorem is formalized as stated, with the same witness, as
\texttt{violatesLowerBound\_152} in
\texttt{lean/Taeyoung/Methods/Negative/Atlas152.lean}: the five masses
\eqref{eq:masses} over the common denominator $3750000$, the cell values
\eqref{eq:cells} over $16$, and the fractional fibre $W=1/16$ between the bulk
class and the first exceptional class.  Both \eqref{eq:edge-density} and
\eqref{eq:hom-density} are evaluated exactly, by kernel reduction over the
$5^6$ assignments, and the comparison with \eqref{eq:target} is exact rational
arithmetic.  There is no difference of substance in the verification.

Two things in the note are not formalized, because the classification does not
use them.  The exact local mechanism of Section~3 --- the kernel matrix
\eqref{eq:negative-kernel-matrix} and the negative quadratic form
\eqref{eq:negative-quadratic} --- explains how the witness was found and plays
no part in checking it, exactly as that section says.  And
\eqref{eq:chromatic-polynomial} is obtained not by deletion--contraction but
from the six counts of surjective colourings, $0,0,0,18,264,840,720$, checked
by kernel reduction; the \texttt{ViolatesLowerBound} statement has to exhibit
the polynomial, so it is proved rather than quoted.

\end{document}
